\section{Conclusion and future directions}
\label{sec:conclusion}

We presented a general in-situ and simultaneous magic state injection method for CSS qLDPC codes. Our scheme prepares multiple logical magic states directly within a single qLDPC block using only standard syndrome extraction resources, without auxiliary preparation blocks for state transfer. This work establishes a qLDPC-native injection primitive for encoded non-Clifford resource states. We showed that the same framework can be substantially improved by optimizing the initial configuration to maximize first-round protection of the injected logical supports. The main advantage of our scheme is the qubit footprint: the BB injection uses 288 physical qubits, compared with 968 for four distance-11 surface code patches. This substantially smaller footprint makes the BB approach relevant for footprint-constrained demonstrations, especially for early fault-tolerant architectures. Circuit-level evaluation on BB and HGP code instances demonstrates both practical operating points on selected BB configurations and general applicability beyond a single code family. 

Several directions remain open. First, increasing the number of simultaneously injectable logical qubits while retaining low error rates will likely require stronger combinatorial conditions, better logical representatives, and improved syndrome extraction circuits. Second, reducing the discard rate remains important for practical throughput; our results suggest that this will require code- and circuit-level improvements, especially in the structure of the round-parity detectors. Finally, a hardware-aware version of the scheme should adapt Bell state preparation on overlapping qubits to the native connectivity, gate set, and resource constraints of a given qLDPC memory experiment. We hope these directions will lead to fully qLDPC-native protocols for high-quality non-Clifford resource state preparation.
